% How a citation prints, shared by the report and its timelines so the two
% cannot differ. The rules themselves are in docs/repository.md, "The
% bibliography's style sheet".

% A footnote names the source and the page; the web address and the access date
% appear only in the bibliography.
% How a law is dated: a case by its year, whatever day it was decided; an act or
% a constitution not at all, since its title already carries the date.
\newcommand{\legaldates}{%
  \ifentrytype{jurisdiction}{\clearfield{month}\clearfield{day}}{}%
  \ifkeyword{datedintitle}{\clearfield{date}\clearfield{year}\clearfield{month}\clearfield{day}}{}}

\AtEveryCitekey{%
  \clearfield{url}\clearfield{urldate}\clearfield{urlyear}\clearfield{urlmonth}\clearfield{urlday}%
  \legaldates}

% The bibliography keeps the web address and drops the access date.
\AtEveryBibitem{%
  \clearfield{urldate}\clearfield{urlyear}\clearfield{urlmonth}\clearfield{urlday}%
  \legaldates}

% Primary law is cited in the notes and kept out of the Works Cited; it is listed
% once, with the pages that cite it, in the Legal Authorities after it
% (docs/repository.md). biblatex-chicago skips these types in any list unless told
% not to, and \printworkscited then leaves them out by type. The documents load
% biblatex-chicago with backref=true, which is what supplies the pages.
\ExecuteBibliographyOptions[jurisdiction,legislation]{skipbib=false}
\defbibfilter{constitutions}{type=legislation and subtype=constitution}
\defbibfilter{statutes}{type=legislation and not subtype=constitution}

% "Cited on page 8" after an authority, in its own sentence.
\renewbibmacro*{pageref}{%
  \iflistundef{pageref}{}{%
    \setunit{\addperiod\addspace}%
    \printtext{\ifnumgreater{\value{pageref}}{1}{\bibstring{backrefpages}}{\bibstring{backrefpage}}%
      \ppspace\printlist[pageref][-\value{listtotal}]{pageref}}}}
\DefineBibliographyStrings{english}{backrefpage={Cited on page},backrefpages={Cited on pages}}

% Works Cited and each legal list keep the look they have always had --- an
% unnumbered \section* for Works Cited, an unnumbered \subsection* for each
% legal list --- and add themselves to the contents at section level, so a
% reader can click to any of them the way to a regular section.
\defbibheading{bibintocsection}{\phantomsection\section*{#1}\addcontentsline{toc}{section}{#1}}
\defbibheading{sublistintocsection}{\phantomsection\subsection*{#1}\addcontentsline{toc}{section}{#1}}

% The sources a reader looks up: everything but the law, without the pages.
\newcommand{\printworkscited}[1]{%
  {\renewbibmacro*{pageref}{}\printbibliography[heading=bibintocsection,title={#1},nottype=jurisdiction,nottype=legislation]}}

% The law: cases by name, constitutions and then statutes by date (a
% legislation entry's sorttitle is its date).
\newcommand{\printauthorities}{%
  \clearpage
  \printbibheading[title={Legal Authorities}]%
  \printbibliography[heading=sublistintocsection,title={Cases},type=jurisdiction]%
  \printbibliography[heading=sublistintocsection,title={Constitutions},filter=constitutions]%
  \printbibliography[heading=sublistintocsection,title={Statutes},filter=statutes]}

% A case name is italic, in the notes and in the table.
\DeclareFieldFormat[jurisdiction]{title}{\mkbibemph{#1}}
\DeclareFieldFormat[jurisdiction]{shorttitle}{\mkbibemph{#1}}
